
\section{CERT-CD}\label{sec:algorithm}

\subsection{The algorithm}\label{sec:algo-desc}

CERT-CD maintains a three-valued status for every unordered pair and a
two-valued status for every ordered pair. It streams data in batches. After
each batch it updates the running Gram matrix, recomputes $A_t(i,j)$ for every
undecided pair, and certifies adjacency when the running maximum of $A_t(i,j)$
crosses the pair's threshold. At that moment the pair's conditioning block is
frozen and its two direction processes are opened. Certified arrowheads are
added when a direction process crosses its own threshold.

Budgets are fixed before any data arrive. The total budget splits as
$\alpha = \alpha_A + \alpha_D$; $\alpha_A$ is divided evenly over the
$\binom{d}{2}$ adjacency nulls and $\alpha_D$ evenly over the $d(d-1)$
direction nulls, so the thresholds are
\begin{equation}\label{eq:thresholds}
\lambda_A = \frac{\binom{d}{2}}{\alpha_A},
\qquad
\lambda_D = \frac{d(d-1)}{\alpha_D}.
\end{equation}
Nothing about the split adapts to the data. Which certificates the algorithm
happens to evaluate---it only opens direction processes for certified
pairs---does not enter the bound, so the graph-dependent order in which pairs
are examined is free.

\begin{algorithm}[t]
\caption{CERT-CD}\label{alg:certcd}
\begin{algorithmic}[1]
\Require stream $o_1, o_2, \dots$; budget $\alpha$; order $k$; warm-up $m$
\State $\alpha_A, \alpha_D \gets$ split of $\alpha$;
       $\lambda_A, \lambda_D$ as in~\eqref{eq:thresholds}
\State consume $o_1,\dots,o_m$ to fix standardisation constants and the prior
       scale $g$; \textbf{discard} them
\State $\mathrm{status}(i,j) \gets \textsc{Undecided}$ for all pairs;
       $\mathrm{arrow}(i,j) \gets \textsc{false}$ for all ordered pairs
\State $\bar A(i,j) \gets 0$; $\bar E(i,j) \gets 0$
       \Comment{running maxima}
\For{each batch $B$}
  \State update the running Gram matrix with $B$
  \ForAll{pairs $(i,j)$ with $\mathrm{status}(i,j) = \textsc{Undecided}$}
    \State $A \gets \min_{S \in \Scal_k(i,j)} E^{(S)}(i,j)$
           \Comment{Schur complements of the Gram matrix}
    \State $\bar A(i,j) \gets \max\{\bar A(i,j), A\}$
    \If{$\bar A(i,j) \ge \lambda_A$}
      \State $\mathrm{status}(i,j) \gets \textsc{CertifiedEdge}$
      \State $Z_{ij} \gets \{\,\ell : \mathrm{status}(i,\ell) \text{ or }
             \mathrm{status}(j,\ell) = \textsc{CertifiedEdge} \,\}$
             \Comment{frozen}
      \State open direction processes for $(i,j)$ and $(j,i)$, starting from
             the next observation
    \EndIf
  \EndFor
  \ForAll{open ordered pairs $(i,j)$}
    \State update $\log E(i,j)$ by~\eqref{eq:direction} on the observations
           since the pair was opened
    \State $\bar E(i,j) \gets \max\{\bar E(i,j), E(i,j)\}$
    \If{$\bar E(i,j) \ge \lambda_D$} $\mathrm{arrow}(i,j) \gets \textsc{true}$
    \EndIf
  \EndFor
  \State \textbf{output} the three-valued graph; the analyst may stop here
\EndFor
\end{algorithmic}
\end{algorithm}

The output graph contains only certified edges. An edge carries an arrowhead
at an end where a direction certificate has fired and a circle otherwise, in
the mark notation of \citet{richardson2002ancestral}. \textbf{The absence of
an edge means undecided, not certified absent.}

\subsection{The whole-graph guarantee}\label{sec:theorem}

\begin{theorem}[Whole-graph anytime validity]\label{thm:main}
Suppose Assumptions~\ref{ass:markov} and~\ref{ass:sep} hold, the per-set
processes $E^{(S)}$ are e-processes for their point nulls, and---for the
direction certificates only---the pairs are generated by the linear
non-Gaussian model of Section~\ref{sec:direction} with the supremum
in~\eqref{eq:direction} attained. Then
\begin{equation}\label{eq:mainbound}
\Prob\left(
\begin{array}{c}
\exists\, t \ge 0 : \text{some edge certified by time $t$ is absent} \\[-1pt]
\text{from } \Gstar, \text{ or some arrowhead certified by time $t$
  is wrong}
\end{array}
\right) \;\le\; \alpha .
\end{equation}
In particular, for \emph{any} stopping time $\tau$---including one chosen by
inspecting the output---the graph reported at time $\tau$ satisfies
$\Prob(\text{it contains a false certified edge or arrowhead}) \le \alpha$.
\end{theorem}

\begin{proof}
Index the adjacency nulls by pairs and the direction nulls by ordered pairs;
both families are determined by $d$ and $k$ and are fixed before any
observation. Consider a single adjacency null for a pair $(i,j)$ that is
non-adjacent in $\Gstar$. By Lemma~\ref{lem:certifying}, $P \in
H^{\mathrm{sep}}_{ij}$; by Proposition~\ref{prop:adjacency}, $(A_t(i,j))$ is
an e-process for that null; so by~\eqref{eq:ville} the probability that
$A_t(i,j)$ ever reaches $\lambda_A = \binom{d}{2}/\alpha_A$ is at most
$\alpha_A / \binom{d}{2}$. Since the algorithm certifies the edge exactly on
that event, the probability that this pair is ever falsely certified is at
most $\alpha_A/\binom{d}{2}$.

Similarly, consider an ordered pair $(i,j)$ for which the arrowhead $i \to j$
is absent from $\Gstar$. Then $P \in H^{\mathrm{dir}}_{j \to i}$, which by
Proposition~\ref{prop:direction} (applied conditionally on $\Fcal_T$, see
Appendix~\ref{app:conditioning}) has $(E_t(i,j))$ as an e-process, and the
probability that it ever reaches $\lambda_D$ is at most
$\alpha_D / (d(d-1))$.

The event in~\eqref{eq:mainbound} is the union, over the at most
$\binom{d}{2}$ non-adjacent pairs and at most $d(d-1)$ absent arrowheads, of
these events. A union bound gives
$\binom{d}{2} \cdot \alpha_A/\binom{d}{2} + d(d-1) \cdot \alpha_D/(d(d-1))
= \alpha_A + \alpha_D = \alpha$.

The final statement follows because the event in~\eqref{eq:mainbound}
contains, for every stopping time $\tau$, the event that the graph at time
$\tau$ contains a false certified assertion; the bound is time-uniform, so no
separate argument per stopping time is needed.
\end{proof}

We stress the form of the quantifier. The theorem does not say ``for each
fixed $\tau$, the error at $\tau$ is at most $\alpha$'', which would be a
family of pointwise statements requiring a further union bound to combine.
It says that the probability of the trajectory \emph{ever} containing a false
certified assertion is at most $\alpha$, and every stopping-time statement is
a consequence. This is the difference between a fixed-sample guarantee applied
repeatedly and a time-uniform one, and Section~\ref{sec:exp-calibration}
measures how large that difference is in practice.

\begin{remark}[What the theorem does not cover]
Theorem~\ref{thm:main} bounds the probability of a false \emph{assertion}. It
says nothing about the pairs left undecided, and it does not bound the
probability that a true edge is missed---there is no such bound, and there
cannot be one without a faithfulness-type assumption. It also does not cover
the analyst's use of the output for anything else: an intervention chosen on
the basis of a certified arrowhead inherits the $\alpha$, but a downstream
effect estimate computed from the certified graph does not automatically
inherit anything, and would need its own analysis; the estimand there is the intervention-effect
problem of \citet{maathuis2009estimating}, and combining a certified graph
with a valid effect interval is a genuinely open question.
\end{remark}

\subsection{Complexity}\label{sec:complexity}

Per batch, the Gram-matrix update is $O(d^2)$ per observation. Recomputing
$A_t(i,j)$ for a pair costs $O(|\Scal_k| \cdot k^3)$, dominated by the Schur
complements, and $|\Scal_k| = \sum_{r \le k}\binom{d-2}{r} = O(d^k)$. Summed
over pairs, an update of the whole adjacency layer is $O(d^{k+2} k^3)$, with
no dependence on $n$: all the data enter through the sufficient statistic.
Only undecided pairs are updated, so the cost falls as the run proceeds.

The direction layer is the expensive part. Each open ordered pair requires a
generalised-Gaussian fit at every batch, which is an IRLS loop plus a scalar
optimisation over $\kappa$, and it does not reduce to a sufficient statistic
because the fit is nonlinear in the residuals. The cost is $O(d^2 \cdot
n_{\text{batch}} \cdot |Z| \cdot \text{iterations})$ per batch, and it is the
reason our experiments stop at $d = 11$. This is an implementation limit, not
a limit of the construction: the numerator's parameters may be refitted at any
frequency (refitting less often only costs power, since the numerator needs to
be $\Fcal_{s-1}$-measurable and nothing more), and the denominator's maximised
likelihood admits warm starts across batches. We have not pursued either
optimisation and regard scaling as unproven; see Section~\ref{sec:discussion}.

\subsection{On the multiplicity layer}\label{sec:multiplicity}

Theorem~\ref{thm:main} uses a union bound over a fixed family, and a union
bound is the least powerful admissible way to combine evidence. It is
reasonable to ask why we do not use a sharper procedure. E-value-based
family-wise error control has advanced considerably: \citet{hartog2026fwer}
show that e-closed testing with weighted arithmetic-mean local tests strictly
dominates e-Bonferroni, and their e-Holm procedure achieves this at $O(m \log
m)$ cost. Crucially, the arithmetic-mean local test is valid under
\emph{arbitrary} dependence between the base e-values, so the fact that all
our per-hypothesis processes are functions of the same running Gram matrix is
\emph{not} an obstacle. An earlier draft of this paper claimed it was; that
claim was wrong, and we correct it here.

The actual obstruction is different and, we think, more interesting. Our
certificates fire on the \emph{running maximum} $\bar A_t = \sup_{s \le t}
A_s$ of an e-process, because that is what a monitored run rejects on. A
running maximum is not an e-value. Ville's inequality bounds
$\Prob(\bar A_\infty \ge 1/\alpha)$ by $\alpha$, which makes $\bar A_t$ what
\citet{hartog2026fwer} call a \emph{pseudo} e-value, but $\Exp[\bar A_\tau]$
generally exceeds one, and closed testing takes e-values as inputs, not
pseudo e-values. Feeding running maxima into an e-closed procedure is not
licensed by the theory. \citet{hartog2026fwer} address this case directly and
recover a bound of $\alpha + O(\alpha^2 \log \alpha^{-1})$ for pseudo
e-values, but only under \emph{independence} across hypotheses---and that is
where our shared sufficient statistic does bite.

So the situation is: dependence is fine for e-closed testing with e-values;
running maxima are fine under independence; we have running maxima under
dependence, and no result covers the combination. The union bound over Ville
events is what remains available, and it is exact for what it does. Closing
this gap---a family-wise procedure for dependent pseudo e-values that improves
on Bonferroni---would immediately improve the power of every method in this
paper, and we regard it as the most valuable piece of the surrounding theory
that is missing.

Two weaker improvements are available now and we note them. First, the budget
split need not be uniform: any fixed weights $w_h > 0$ with $\sum_h w_h = 1$
give thresholds $1/(\alpha w_h)$ and the same bound, so prior knowledge about
which pairs are likely adjacent can be spent as unequal weights, decided
before the data arrive. Second, if the analyst wants false-discovery-rate
rather than family-wise control, the e-BH procedure of
\citet{wang2022false} applies to e-values under arbitrary dependence, as do
the martingale-based merging rules of \citet{vovk2024merging}, and would
replace the union bound directly---but only at a fixed stopping time, since it
again takes e-values rather than running maxima as inputs.

\subsection{An optional fourth verdict: certified negligible}
\label{sec:negligible}

Section~\ref{sec:principle} argued that there is no certifying null for
non-adjacency. There is, however, one for a weaker proposition, and analysts
who need it can have it at the cost of a stated margin. Fix $\delta > 0$ and
consider the interval null
\begin{equation}\label{eq:ropenull}
H^{\mathrm{big}}_{ij,S} \;:\quad |\rho_{ij \cdot S}| \ge \delta ,
\end{equation}
where $\rho_{ij\cdot S}$ is the partial correlation. Rejecting it for
\emph{every} $S$ in a chosen set certifies that the pair's association is
smaller than $\delta$ given each of those conditioning sets---a statement
about effect size, not about the graph. The confidence
sequence~\eqref{eq:cs} supplies the test directly: reject when the interval
excludes $[-\delta, \delta]$. The resulting verdict, \emph{certified
negligible at margin $\delta$}, is honest in a way that a deleted edge is not,
because $\delta$ is reported alongside it and the reader can judge whether an
association below $\delta$ matters for their purpose. It is also strictly less
than the claim a delete-on-non-rejection method makes, which is presumably why
that claim is popular. We implement this verdict but do not use it in the
experiments below, since the comparisons there are with methods that assert
absence outright.

\section{Latent confounders}\label{sec:latent}

The adjacency certificate does not require causal sufficiency; nothing in
Lemma~\ref{lem:certifying} assumes that all common causes are observed.
What changes under latent confounding is the interpretation of a certified
edge and the availability of the direction certificate.

If $\Gstar$ is a DAG over observed and latent variables and we observe only
$V$, the relevant object is the maximal ancestral graph over $V$
\citep{richardson2002ancestral}, and the relevant separation notion is
m-separation. Assumption~\ref{ass:sep} becomes the requirement that every
non-adjacent pair in the MAG has an m-separating set of size at most $k$ among
the observed variables---which, unlike the DAG case, is not guaranteed by
bounding an in-degree, and is why FCI searches Possible-D-SEP sets rather than
neighbourhoods. With that caveat, Lemma~\ref{lem:certifying} and
Proposition~\ref{prop:adjacency} go through verbatim with d-separation
replaced by m-separation, and a certified edge means ``$i$ and $j$ are
adjacent in the MAG'', that is, neither is separable from the other by any
observed set---which is compatible with a direct effect, a latent common
cause, or both.

We therefore implement CERT-FCI: run the adjacency certificates to obtain a
certified skeleton, apply the FCI orientation rules R1--R4 \citep[originally
][in the formulation of \citealp{zhang2008completeness}]{spirtes1995causal} to the certified edges only, and
report the resulting partial ancestral graph. Orientations produced by the FCI
rules inherit the adjacency certificates' guarantee only in the weak sense
that they are sound given a correct skeleton and correct sepsets; we do not
claim a family-wise bound for them, and we flag this as the principal gap in
the latent-variable extension. The direction certificate of
Section~\ref{sec:direction} is unavailable here, because
Proposition~\ref{prop:lingam} assumes no unobserved common cause of the pair.
Section~\ref{sec:exp-latent} reports what CERT-FCI recovers, and
Section~\ref{sec:exp-violations} reports what latent confounding does to the
direction certificate when it is applied anyway.
